% Technical Background, section 4. Pointed back to by: methodology/05_curvature.tex.
% Model: the derivation chain of Aida 3.3.1 (guide section 5.1), name the assumption at each step.
% A single section without subsections; this is the chapter's main equation chain.
% Cite: zhang2024curvature, telloayala2026tortuosity (only as a discrete curvature estimator, the
%   turning angle; its CAD results belong in the literature review).
% CITE WANTED: a differential geometry of curves source (e.g. do Carmo).
% Split with the Methodology: kappa(s), and the turning angle over equal steps as its discrete
%   estimate (atan2 form), are defined HERE. The Methodology keeps only what carries a setting: the
%   per-artery path, the 5 mm chord, T_ell = sum(theta)/L_c and the ordered angle sequence.
% No hazard closer: the noise problem is stated as fact in P4 and resolved in the Methodology.

\section{Curvature of a Centerline}
\label{sec:tech-curvature}
% P1 definition opener: a centerline as a space curve c(s) parametrised by arc length.
% P2 curvature kappa(s) = ||c''(s)||, the inverse radius of the osculating circle; equation + where.
%    It does not depend on the position or orientation of the curve, so it needs no common frame.
%    Draft prose below moved from clinical_background_v2/04_flow_to_geometry.tex (2026-10-05); add
%    the equation when P1 is written.

Curvature is defined at a point on the centerline as the inverse of the radius of the circle that
best fits the curve there, so that a tight bend has high curvature and a straight segment has none.
The mean curvature is its average along the vessel, whereas the curvature profile is the curve of
curvature against distance from the origin of the vessel, preserving where along the vessel a bend
is located.
The tortuosity index, the ratio of a vessel's length to the straight-line distance between its
ends, instead summarizes the whole vessel in one number and cannot distinguish a vessel with one
sharp bend from one with several gentle bends (Figure~\ref{fig:curvature}).

\begin{figure}[h]
  \centering
  \includegraphics[width=0.85\textwidth]{content/background/figures/rq_curvature.pdf}
  \caption{The tortuosity index cannot tell these three vessels apart, but curvature can.
    Each synthetic centerline is 50\,mm long between ends 40\,mm apart, so all three share the
    tortuosity index of 1.25, yet their mean curvature differs sixfold. Color shows local
    curvature $\kappa$.}
  \label{fig:curvature}
\end{figure}

% P3 a centerline is a sequence of points, so derivatives are estimated: finite differences, fitted
%    splines, or the turning angle between consecutive segments (telloayala2026tortuosity).
% P4 each derivative amplifies point noise, the second more than the first, so curvature is
%    estimated after smoothing or over a longer step, and depends on that choice; for the turning
%    angle, false turning from point noise sigma scales as sigma / step^2 per unit length.
